\section{Conclusions}
\label{vii:sec:conclusiones}

Torsional amplitude is determined by the spinorial action and material
state, with an explicit formula and a sector in which its positivity is
preserved. This result yields an energetic contribution proportional to
\(-a^{-6}\), a unique homogeneous bounce under the proved conditions,
and, for flat dust with \(\Lambda_{\mathrm{eff}}=0\), a regular evolution
for all time with complete Levi--Civita causal geodesics. The sequence
connects a quartic state observable to a minimum cosmological scale. Its
origin remains in the APP--TRIT--TPK construction: orientation,
remainder, carry, and depth are transported through to the actions,
currents, and observables of the geometric realization.

\subsection{Determination of the amplitude and cosmological evolution}

The central result of the spinorial realization is the determination
of torsional amplitude as a functional of the material state:
\[
 s_0[\varrho,V_c]
 =\frac{3\kappa c^2\hbar^2}{16}
   \frac{\gamma^2}{1+\gamma^2}
   \frac{\operatorname{Tr}
    (\varrho\widehat{\mathscr Q}_{\rm sp})}{V_c^2}.
\]
The formula follows successively from the internal Clifford
representation, the spin current, elimination of contorsion,
the normally ordered quartic observable, and lapse and
volume variations. The quadratic volume factor expresses the
dependence of density on the second moment of the
current. A zero mean spin retains a second moment
that can contribute to geometry.

In the occupation-three sector one obtains
\(\widehat{\mathscr Q}_{\rm sp}=4I\). For trace-one states supported
in that sector, every unitary evolution preserving the sector preserves
its expectation and determines a positive amplitude.
The material state and comoving cell individualize the realization;
\(G\), \(\hbar\), and \(\gamma\) retain their structural
antecedents. The reduced Planck constant enters quadratically into
the interaction, Barbero--Immirzi determines the coefficient of
the Holst operator, and gravitation converts that interaction into
a geometric response. This composition jointly explains
the functional value of the amplitude and the role of each factor.

The memory law \(M_n=\varphi^{-2n}\) and spatial reading
\(a_n/a_{\rm ref}=\varphi^{n/3}\) produce dilution
\(a^{-6}\). In the spinorial realization with preserved second
moment, the same power appears through material variation.
With \(\rho_0,s_0>0\), \(w<1\), zero spatial curvature, and
\(\Lambda_{\rm eff}=0\), the bounce threshold is unique and
transverse. Two distinct bounds locally preserve its
existence, uniqueness, and time orientation under perturbations.

The dust sector also admits an exact solution for all
time. The scale factor remains bounded away from zero,
curvature is finite, and the causal Levi--Civita geodesics
are complete in both directions. The global result belongs
to that explicit realization; persistence of the bounce is
established under the local conditions of its theorem. Both
proofs show how the amplitude obtained from the state translates
into a minimum scale and an evolution through contraction and expansion.

\subsection{From transported memory to gravitational response}

Radial normalization follows the longitudinal composition
\(L_\alpha=(5759/23040)\alpha^{16}L_*\). Normalized archive transport
and the positive metric yield
\(R_X=L_\alpha\mathsf U X\mathsf U^*\); energy preserves the same
character and its conjugate length satisfies
\(H_X\Lambda_X=\hbar_{\rm ret}cI\). Under the reduced-radius reading,
these identities determine \(G=c^3L_\alpha^2/\hbar_{\rm ret}\), common
to the positive modes of the realization. The proof retains its
longitudinal and metric rules and the stated physical identification.
Joint boundary and memory reduction preserves the coefficient; dynamic
reduction also transports the spectral parameter and the temporal norm.
The circular expression is recovered with
\(t_0=\pi_{\rm HMT}L_\alpha/(54c)\).

The five-component space is obtained through an explicit projection
of the twelve vertices and is transported isometrically to symmetric
trace-free tensors. Transverse reduction preserves two components.
The sequence \(12\to5\to5\to2\) distinguishes the initial incidence,
the gravitational carrier, its tensor representation, and the transverse
reading. The coefficient \(G_a\) fixes the common coupling;
orientations belong to the operator and its components.

In the same return section, the spinorial-amplitude coefficient has the
equivalent expression
\[
 s_0[\varrho,V_c]
 =\frac{3\pi c\hbar_{\rm ret}L_\alpha^2}{2}
   \frac{\gamma^2}{1+\gamma^2}
   \frac{\operatorname{Tr}(\varrho\widehat{\mathscr Q}_{\rm sp})}{V_c^2}.
\]
It follows by substituting
\(\kappa=8\pi L_\alpha^2/(\hbar_{\rm ret}c)\) into the proved
functional. Length and action therefore fix its coefficient, while the
state and volume retain their material roles in the cosmological solution.

The memory screen provides soldering and an energetic response
compatible with refinement. Its minimizing extension,
constructed from transported incidence, determines an
energy and a differential with respect to the connection. The calculation
includes variations of incidence, boundary, weight, and normalization.
Upon passage to the Cartan realization, the co-tetrad and the material
pairing convert this differential into a current with an explicit
geometric domain.

The connection equations admit algebraic inversion of the
Holst and torsion operators. In the material sector affine in
contorsion, this inversion produces the effective correction to
the action. For the minimal extension, whose current generally
depends on the connection itself, the complete stationarity equation,
the local continuation criterion through its Hessian, and the
integral formula for the effective action are proved. The spinorial and
minimizing-extension actions therefore retain their own arguments and
variation laws. The former supplies the amplitude functional evaluated
in the cosmological solution; the latter describes the full variational
response of the transported screen.

\subsection{Information, horizons, and relational response}

The residual balances bring together the trace, the
trace-free part, and their divergences in one formulation. With constant
couplings in the chart and a conserved material source, the residual tensor
covariantly preserves its total balance; variation of its
trace determines exchange between the two sectors. A constant trace
characterizes separate conservation.

Horizon information uses the areal unit,
action, and thermal conversion previously constructed.
The reduced state and its purification distinguish accessible
information from preserved correlations. In the spherical balance
developed here, \(TS=E\),
\(T\,dS=2\,dE\), and \(S\,dT=-dE\) are jointly verified;
the dependence of temperature on radius is part of these identities.
The Boltzmann constant converts informational content
within the thermal section, while the action--gravitation relation
determines the area unit of its geometric reading.

\Needspace{24\baselineskip}
The areal and chronological composition determines
\(\mathcal T_{I/A}q_{\mathcal A}/\log3=h/(108t_0\log3)\).
At \(R=L_\alpha\), cosmological and chronological temperatures
are related by the factor \(\log3/(2\pi)\).
For the response Hamiltonian
\(\mathsf H_R=\tau_R\Delta_0\mathsf D_\partial\), the weights
90 and 120 correspond to different energy gaps with a common scale.
The finite balance is exactly
\(\mathcal E_R-T_{\rm cos}\mathcal S_R
=\tau_R\mathcal D(\sigma\Vert\rho_0)\), and its tangent term
is the first law at fixed radius. The relative section constructed from
that state determines a unique positive transducer; identification of
an additional reader reduces to checking its value on the chronological
section in~\eqref{viii:eq:criterio-una-seccion}. Record-preserving
isometries transport this balance and its observables.

The observer is defined through operations of orientation,
trivialization, and restriction of observables. The self-inertial
response preserves its factorization through global information
and makes the defect produced by projection explicit.
Finally, in the realized chronological channel, the covariant difference
constructed with route transport admits retarded and advanced
Green operators for finitely supported sources. Compatibility of both readings
at the boundary specifies a source space and a projection
entering the variational response. Bidirectionality
is expressed by operators, moments, and boundary conditions
that can be checked within the article.

The decisive relationship is that information preserved by the state
enters the geometric response through a current and its second moment.
Elimination of contorsion and metric variation convert that observable
into a density and pressure; the dust solution shows how both determine
a minimum scale and a causally complete evolution. Balances and the
observer's reading preserve, in turn, the distinction between accessible
information and global correlation. Within the constructed domains,
memory, material source, and geometry thus form an explicit chain of
determinations, from the arithmetic kernel to gravitational evolution
and its informational content.
